\chapter{Conclusions and Future Work}
\label{chap:conclusions}

\section{Achievements}
\label{sec:conc-achievements}

The work documented in this thesis delivered a Network Intrusion
Detection System that meets the three engineering questions set
out in the introduction.

The first question --- whether signature and machine-learning
detection can cooperate explicitly in a single application ---
is answered by the dispatcher's twenty-line merge logic and the
three alert categories it produces. The synthetic validation in
Chapter~\ref{chap:validation} shows that the hybrid category
forms approximately two thirds of the alerts in the test
catalogue once a Random Forest is activated. The pattern is
exactly the cooperative behaviour the project set out to
exhibit: signatures contribute precision and explanation, ML
contributes corroboration, and the merge logic turns the
combination into a single user-visible label.

The second question --- whether per-alert forensic context can
be persisted without an external packet sniffer --- is answered
by the JSONL ring-tail mechanism. Every alert in the database
has a path to its capture file, every capture file contains the
last fifty packets the dispatcher saw, and the file is
downloadable from the UI through an authenticated endpoint. The
synthetic validation in Chapter~\ref{chap:validation} confirms
that the mechanism produces a complete capture for each alert
fired, with source and destination IPs preserved exactly as
they entered the dispatcher.

The third question --- whether the system can be deployed without
infrastructure expertise --- is answered by the
\file{run.py} launcher and its two wrappers. A user with Python
3.12 and Node.js installed can clone the repository, double-click
\file{run.bat}, and have a working SIEM in their browser within
two minutes. The launcher handles the venv creation, the
dependency installation, the subprocess supervision and the
shutdown cleanly. The reliability work documented in
Section~\ref{sec:design-launcher} (idempotent shutdown, Python
version check, dependency stamp file, fallback to backend-only
mode) is the difference between a thesis prototype and software
that other people can actually run.

Beyond the three originally stated questions, the project
acquired a fourth contribution during development: the integrated
\emph{Attack Console}. The component was added to solve the
demo problem of needing on-demand attacks during a live
presentation and turned into a small-scale, deterministic test
harness for the whole pipeline. The Console exposes six attack
types as UI buttons, fires them through the same dispatcher path
that a real network interface would feed, and reports the number
of alerts generated as immediate visual feedback. It became the
most demonstrable feature of the entire system and is, to my
knowledge, not present in any of the related-work tools listed
in Chapter~\ref{chap:related-work}.

\section{Limitations}
\label{sec:conc-limitations}

Honest accounting requires listing what the project does
\emph{not} do. The most significant limitations, in rough order
of importance, are:

\textbf{The ML detector has only been evaluated on synthetic
data.} The training set is generated inside
\file{ml\_detector.py} from distributions designed to mimic
benign and malicious traffic. The model performs almost
perfectly on its own held-out test split, but that says little
about how it would perform on real attack traffic outside the
catalogue. The natural remediation is evaluation on
CICIDS2017~\cite{sharafaldin2018cicids} or
NSL-KDD~\cite{tavallaee2009detailed}, both of which are widely
used in the IDS literature for exactly this purpose.

\textbf{Per-packet features rather than per-flow.} The six
features in Section~\ref{sec:design-features} describe a
single packet. Many real-world attack patterns (slow scans,
data exfiltration, beaconing) are only visible at the flow level
--- packets aggregated into a bidirectional 5-tuple session ---
and the current feature vector cannot capture them. A
per-flow refactor is the largest single piece of remaining work
and would be the first thing I would do if I had another
semester on the project.

\textbf{Schema evolution via \code{create\_all}.} The database
schema is created at startup by SQLAlchemy's
\code{Base.metadata.create\_all()}. This works for the initial
deployment but loses data on schema changes; a real deployment
would use Alembic migrations. The work to add Alembic is
mostly mechanical, but I did not consider it justified within
the thesis timeline.

\textbf{No alert correlation.} Each alert is independent. In a
production environment a single Nmap scan would produce dozens
of alerts that should be grouped into a single ``incident'';
the current console shows them as separate rows. Wazuh and
Splunk Enterprise Security both have correlation rule engines;
\nidsname{} has none.

\textbf{Authentication is single-role.} The seeded admin user
can do everything. There is no read-only analyst role, no
audit log of who edited which rule, no two-factor
authentication. For thesis-scale demonstration these are
acceptable omissions; for a real deployment they are not.

\textbf{Live capture untested on macOS.} I have run the
LiveSniffer on Linux (libpcap) and Windows (Npcap) but not on
macOS. The path should work because Scapy supports macOS, but
I cannot personally certify it without testing.

\section{Future Work}
\label{sec:conc-future}

The limitations above translate directly into a roadmap. In rough
order of priority:

\begin{enumerate}
  \item \textbf{Flow-level features.} Aggregate packets into
    bidirectional 5-tuple sessions, compute flow statistics
    (packet count, byte count, mean and standard deviation of
    inter-arrival, protocol-specific fields), and refit the ML
    detector on the richer feature vector. Expected to improve
    detection of multi-packet attack patterns.
  \item \textbf{Evaluation on public datasets.} Train and
    evaluate Random Forest and Isolation Forest on
    CICIDS2017, report standard metrics (precision, recall,
    F1, AUC) per attack family, and add a chapter to the
    thesis comparing the results against published baselines.
  \item \textbf{End-to-end live-capture validation.} Drive the
    same attack catalogue from an external attacker host
    (e.g.\ a Kali Linux virtual machine on the local LAN)
    against the live-capture path, demonstrating that the
    detection numbers reported in Chapter~\ref{chap:validation}
    transfer to genuine network traffic without any code
    change.
  \item \textbf{Alembic migration suite.} Replace
    \code{create\_all()} with a versioned migration system.
    Write the initial migration that represents the current
    schema as the baseline, so that future thesis iterations
    can evolve the schema without losing the demo database.
  \item \textbf{Alert correlation engine.} Implement an
    ``incident'' table that groups related alerts by source IP
    and category within a configurable time window. Surface
    incidents as the default unit of the Alerts page, with
    individual alerts accessible by drill-down.
  \item \textbf{Snort rule import.} Translate a subset of the
    open Snort Emerging Threats rule set into the project's
    JSON match specification, mostly to enrich the default
    rule database and validate the JSON schema's
    expressiveness against a battle-tested ecosystem.
  \item \textbf{User roles and audit log.} Add an
    \code{analyst} role that has read-only access to alerts
    and a write-restricted view of rules. Log every rule
    modification, ML training and forensic download in an
    \code{audit\_log} table.
\end{enumerate}

A few exploratory directions outside the priority queue, but
worth recording:

\begin{itemize}
  \item \textbf{LLM-assisted triage.} A chat panel on the
    Alert Detail page that summarises the captured packets in
    natural language and suggests likely attack categories.
    This is a separate research-grade question and I would only
    pursue it after the priority list is complete.
  \item \textbf{Distributed sensors.} Multiple capture
    instances reporting to a single dispatcher. Useful for
    multi-segment networks but a significant architectural
    change.
  \item \textbf{Production storage backend.} Swap SQLite for
    PostgreSQL or DuckDB. The SQLAlchemy abstraction makes
    the change largely mechanical, but the operational story
    becomes more complex.
\end{itemize}

\section{Closing Remarks}
\label{sec:conc-closing}

When I started this thesis I expected the hardest part to be
the machine learning. It was not. The ML implementation took
roughly a week of work; the harder problems turned out to be the
ones I had not anticipated: getting Scapy to behave on Windows,
debugging the WebSocket auth handshake, writing a launcher that
gracefully handles \verb|Ctrl+C| on three different operating
systems, and making the React WebSocket subscriptions survive
page navigation. None of these are research problems --- they
are engineering problems, and the time spent on them is the
difference between a paper sketch and a system that actually
runs.

The reward for that work is a project that demonstrates the
detection pipeline end-to-end, with a console an analyst can
actually use, deployable with a single double-click. That is
what \nidsname{} is, and I think it is a defensible answer to
the original brief.
